\chapter{Crystal Meth}

\section*{Definition and Overview}
Crystal methamphetamine, commonly called crystal meth or ice, is a highly addictive stimulant drug. It is a form of methamphetamine, a synthetic drug that strongly stimulates the central nervous system, producing intense euphoria, energy, and alertness. Crystal meth is smoked, snorted, injected, or swallowed, and its effects come on quickly and last for many hours. It is associated with serious physical and mental health harms, including psychosis, aggression, heart problems, weight loss, and addiction. Treatment is available and effective, including counselling, support services, and, in some cases, medication for related problems.

\section*{Epidemiology}
Methamphetamine use is a significant public health issue in Australia, with crystal meth the most common form used. Rates of use are highest among young adults, and use is more common in men. A small but important proportion of people who use crystal meth develop dependence and serious harms, including psychosis, violence, and overdose. The drug is linked to substantial emergency department presentations, mental health admissions, and crime. Rates of use have fluctuated, but crystal meth remains one of the most harmful illicit drugs in Australia.

\section*{Aetiology and Pathophysiology}
Crystal meth acts powerfully on the brain's reward and alertness systems.

\begin{itemize}
    \item \textbf{Dopamine release:} methamphetamine causes a large release of dopamine, producing euphoria and motivation
    \item \textbf{Noradrenaline release:} increases alertness, heart rate, and blood pressure
    \item \textbf{Long-lasting effects:} the drug's effects last many hours, and repeated use depletes dopamine stores
    \item \textbf{Addiction:} strong reinforcement leads to compulsive use and dependence
    \item \textbf{Tolerance and withdrawal:} regular use requires more drug for the same effect, and stopping causes cravings, fatigue, and depression
    \item \textbf{Toxic effects:} high doses cause overheating, seizures, heart damage, and stroke
    \item \textbf{Neurotoxicity:} chronic use can damage brain cells involved in mood and memory
\end{itemize}

\section*{Clinical Features}
The effects and harms of crystal meth are wide-ranging.

\begin{itemize}
    \item Intense euphoria, increased energy, and talkativeness shortly after use
    \item Reduced appetite, weight loss, and poor sleep
    \item Raised heart rate, blood pressure, and body temperature
    \item Sweating, dilated pupils, and jaw clenching
    \item Anxiety, agitation, paranoia, and aggression
    \item Psychosis: hallucinations and delusions, which can persist
    \item Skin picking, dental decay ("meth mouth"), and poor self-care
    \item Financial, relationship, and work problems
    \item Withdrawal: exhaustion, depression, and strong cravings when stopping
\end{itemize}

\section*{Differential Diagnosis}
The harms of crystal meth can resemble other conditions.

\begin{itemize}
    \item Other stimulant use, such as cocaine or ecstasy
    \item Mental illness, including schizophrenia, bipolar disorder, and severe anxiety
    \item Drug-induced psychosis from other substances
    \item Hyperthyroidism, with similar agitation and weight loss
    \item Sleep deprivation and burnout
    \item Heart disease or stroke, which can complicate stimulant use
\end{itemize}

\section*{When to Seek Medical Care}
Seek help early if you or someone you care for is using crystal meth. GPs, drug and alcohol services, and counselling are available, and treatment works best when started early.

\textbf{Call 000 (or local emergency services) for:}
\begin{itemize}
    \item Chest pain, difficulty breathing, or a seizure after using
    \item Very high temperature, confusion, or collapse
    \item Severe agitation, violence, or threatening behaviour
    \item Hallucinations or paranoia that are distressing or dangerous
    \item Suicidal thoughts or self-harm
    \item A suspected overdose
\end{itemize}

\section*{Diagnosis and Investigations}
Assessment identifies the extent of use and any related health problems.

\begin{itemize}
    \item \textbf{History:} patterns of use, route, frequency, and impact on life
    \item \textbf{Mental health assessment:} screening for psychosis, depression, anxiety, and suicide risk
    \item \textbf{Physical examination:} weight, teeth, skin, and signs of injection
    \item \textbf{Blood pressure and heart checks:} as stimulants affect the cardiovascular system
    \item \textbf{Blood and urine tests:} for related infections, including hepatitis and HIV
    \item \textbf{Withdrawal assessment:} planning for safe management of cravings and depression
\end{itemize}

\section*{Management}
Treatment is available and involves medical, psychological, and social support.

\begin{itemize}
    \item \textbf{Withdrawal management:} supervised support to stop using safely, managing cravings, sleep, and mood
    \item \textbf{Counselling:} cognitive behavioural therapy and motivational interviewing
    \item \textbf{Residential rehabilitation:} structured programs for people with severe dependence
    \item \textbf{Peer support:} groups such as SMART Recovery
    \item \textbf{Treatment of mental health problems:} care for psychosis, depression, and anxiety
    \item \textbf{Medical care:} treatment of heart, dental, skin, and infectious complications
    \item \textbf{Social support:} housing, employment, and family support
    \item \textbf{Harm reduction:} information on safer use and overdose prevention
\end{itemize}

\section*{Complications}
\begin{itemize}
    \item Psychosis, which can be prolonged and needs treatment
    \item Heart attack, stroke, and seizures
    \item Overheating and dehydration
    \item Severe weight loss and malnutrition
    \item Dental decay and gum disease
    \item Skin infections and abscesses from injecting
    \item Hepatitis, HIV, and other infections from sharing equipment
    \item Overdose and death
    \item Relationship breakdown, job loss, and financial ruin
    \item Injury to the unborn child in pregnancy
\end{itemize}

\section*{Prognosis and Outcomes}
Recovery from crystal meth dependence is achievable, and many people do well with treatment and support. Outcomes are better when dependence is recognised early and when treatment addresses mental health, social circumstances, and relapse prevention together. The intense withdrawal period lasts weeks, and relapse is common, but most people who persist with treatment achieve long periods of abstinence. Even for people with methamphetamine psychosis, most recover with abstinence and treatment, though some develop persistent mental illness.

\section*{Prevention and Self-Care}
\begin{itemize}
    \item Avoid experimenting with the drug; dependence can develop quickly
    \item Seek help from a GP or drug and alcohol service early
    \item Talk openly with trusted people and access counselling
    \item Develop healthy routines: sleep, exercise, nutrition, and social connection
    \item If injecting, use sterile equipment and never share
    \item Have a plan for emergencies and know the overdose signs
    \item Access support services for housing, work, and family
    \item Seek help for any suicidal thoughts immediately
\end{itemize}

\section*{Sources and Further Reading}
The following reputable sources provide further information for healthcare students and patients seeking reliable, current advice about crystal methamphetamine.

\begin{itemize}
    \item Alcohol and Drug Foundation. \textit{Crystal methamphetamine (ice).} https://adf.org.au/
    \item Healthdirect Australia. \textit{Crystal methamphetamine (ice).} https://www.healthdirect.gov.au/
    \item Australian Government Department of Health and Aged Care. \textit{Methamphetamine treatment.} https://www.health.gov.au/
    \item Cracks in the Ice. \textit{Evidence-based information about ice.} https://cracksintheice.org.au/
    \item DirectLine. \textit{Alcohol and drug counselling and referral.} https://www.directline.org.au/
    \item NHS. \textit{Methamphetamine: effects and help.} https://www.nhs.uk/
\end{itemize}
